\documentclass[10pt,a4paper]{amsart}
\usepackage[margin=19mm]{geometry}
\usepackage{amsmath,amssymb,mathtools}
\usepackage[scr=rsfs,cal=euler]{mathalfa}
\usepackage{xfrac}
\usepackage[unicode,hidelinks,bookmarks=false]{hyperref}
\newcommand{\ie}{i.e.\ }
\newcommand{\cf}{cf.\ }
\newcommand{\resp}{respectively\ }
\newcommand{\Subsection}{\subsection}
\providecommand{\nobreakdash}{\nobreak\hbox{-}}
\setlength{\emergencystretch}{2em}
\multlinegap0pt
\usepackage{bm}
\usepackage{bbm}
\usepackage{dsfont}
\usepackage{xspace}
\usepackage{cancel}
\usepackage{mathtools}
\usepackage{amsthm}
\makeatletter
\@ifundefined{thm}{\newtheorem{thm}{Thm}}{}
\@ifundefined{lem}{\newtheorem{lem}[thm]{Lemma}}{}
\makeatother
\newcommand{\QQ}{{\mathbb Q}}
\newcommand{\ZZ}{{\mathbb Z}}
\title[Explicit Families of Hyperelliptic Curves with CM Jacobians]{Explicit Families of Hyperelliptic Curves with CM Jacobians}
\author{Saeed Tafazolian, Jaap Top}
\date{Source: arXiv:2511.08422v1 (CC BY 4.0)}
\begin{document}
\maketitle

\noindent\emph{Source and licence.} Explicit Families of Hyperelliptic Curves with CM Jacobians. Version arXiv:2511.08422v1 (CC BY 4.0), distributed under the Creative Commons Attribution 4.0 International licence, as recorded in the arXiv licence metadata for this exact version and shown on the abstract page https://arxiv.org/abs/2511.08422v1. Full rights notice: licenses/arxiv-2511.08422-CC-BY-4.0.txt.\par\smallskip

\noindent\textit{Editorial scope.} Counted unit: the Lem whose source label is \texttt{fields} in the article (source0.tex lines 268--282, source label \texttt{fields}), delivered together with the article's own complete original proof of it (source0.tex lines 283--323, one \texttt{proof} environment of 41 lines). No mathematics is written, repaired or re-derived: statement and proof are the article's own text line for line. Required context retained uncounted: none. Every definition, display and label of those lines is retained unchanged. Omitted from this package and not redistributed: 1--267, 324--854. Nothing inside a retained excerpt is omitted or altered: every retained body is delivered exactly as the article's lines are written, so no omission receipt of any kind accompanies this package. Citations: the retained excerpts carry no citation command at all, so no bibliography entry of the article is transcribed and no reference is redistributed. Reference adaptation: none was needed and none was made; every reference inside the delivered unit resolves inside it, so no unresolved reference or citation is produced. Printed slips kept literal: no printed word, symbol, hypothesis or number of the counted statement or of its proof was altered. Layout and class: the article's own class is \texttt{amsart} with packages including graphicx, amsmath, amsfonts, amssymb, amsthm, mathrsfs, color, url; the wrapper is the lane's pinned \texttt{amsart} editorial wrapper at 10pt on A4, which restates the theorem-like environments the retained text needs and adds the article's own macro definitions that the retained text needs (\texttt{\textbackslash QQ}, \texttt{\textbackslash ZZ}), with extra wrapper packages (bm, bbm, dsfont, xspace, cancel, mathtools). The delivered numbering is the unit's own. Assets: no retained span contains a figure environment, a graphics-inclusion command, a PDF-page inclusion, a table or a TikZ picture; no image file is copied into this package, the package declares no asset dependency and no image file is redistributed with it. Licence: version arXiv:2511.08422v1 of this article is distributed under the Creative Commons Attribution 4.0 International licence, as recorded in the arXiv licence metadata for this exact version and shown on the abstract page https://arxiv.org/abs/2511.08422v1. A full rights notice is delivered at licenses/arxiv-2511.08422-CC-BY-4.0.txt. Characters: every retained excerpt is pure ASCII, so the delivered engine, XeLaTeX with its default Latin Modern text font, reports no missing character.
\begin{lem}\label{fields}
    \begin{enumerate}
        \item $K_d$ is a CM-field precisely when $d\geq 3$.
        \item In case $4\nmid d$ one has $K_d=\QQ(\zeta_d)$.
        \item If $4|d$, then the CM-field $K_d\subset\QQ(\zeta_d)$ by the
        Galois correspondence relates to the subgroup
        of $(\ZZ/d\ZZ)^\times$ generated by
        $\overline{-1+d/2}$.
        \item The subgroup mentioned in (3) is trivial when
        $d=4$ and has order $2$ otherwise.
        \item In case $d=2^e\geq 8$, the Galois group of
        the CM-field $K_d$ over $\QQ$ is a cyclic group of order $2^{e-2}=d/4$.
        \item Every proper subfield of $K_{2^e}$ is totally real.
    \end{enumerate}
\end{lem}

\begin{proof}
  Assertion (1) follows from the discussion preceding the
  statement and the observation that $i\sin(2\pi/d)$ is real
  only when $d=1,2$.\\
  To prove (2) and (3) and (4), one observes that an automorphism of $\QQ(\zeta_d)$
  coming from $\overline{n}\in (\ZZ/d\ZZ)^\times$ fixes $\zeta_d-\zeta_d^{-1}$
  precisely when $\sin(2\pi/d)=\sin(2\pi n/d)$. The latter equality means
  that one of 
  \[\begin{array}{ccl} (\textrm{a}) &&1\equiv n\bmod d,\\ (\textrm{b})&&2\equiv d-2n\bmod 2d\end{array} \]
  holds. Case (a) means that we have the identity automorphism.
  Case~(b) can only occur when $2|d$, and then $n\equiv -1+d/2\bmod d$. However,
  as $n$ is supposed to be a unit modulo $d$, this also means that $d/2$ has
  to be even, i.e., $4|d$. In conclusion, the subgroup of $(\ZZ/d\ZZ)^\times$
  related to $\QQ(\zeta_d-\zeta_d^{-1})\subset\QQ(\zeta_d)$ in the 
  Galois correspondence, is trivial when $4\nmid d$. Hence in those
  situations the fields coincide, proving (2). Assertion (3) is clear from
  the same conclusion, and (4) is simply the observation that
  $-1+4/2=1$ and for $4|d$ such $d>4$ one has $\overline{-1+d/2}\in (\ZZ/d\ZZ)^\times$ has order $2$.

  To prove (5), from (3) and (4) it follows, with $d=2^e\geq 8$, that $\textrm{Gal}(K_d/\QQ)$ is isomorphic to $(\ZZ/d\ZZ)^\times/\langle \overline{-1+d/2}\rangle$.
  For completeness we include the following elementary reasoning. Consider
  the short exact sequence
  \[ 1\to \{\overline{n}\in (\ZZ/2^e\ZZ)^\times\;:\;n\equiv 1\bmod 4\}\hookrightarrow (\ZZ/2^e\ZZ)^\times\stackrel{\bmod 4}{\longrightarrow}
  (\ZZ/4\ZZ)^\times\to 1.\]
  The group on the left has order $2^{e-2}$, and $5\bmod 2^e$ is contained in it.
  Since $5^{2^{e-3}}\equiv 1+2^{e-1}\bmod 2^e$, the order of $5\bmod 2^e$
  equals $2^{e-2}$. Hence $\overline{5}$ generates the group on the left.
  From $\overline{-1+2^{e-1}}\not\in\langle \overline{5}\rangle$ one now
  concludes that the homomorphism $\langle \overline{5}\rangle \to
  (\ZZ/2^e\ZZ)^\times/\langle \overline{-1+2^{e-1}}\rangle$ given by
  $\overline{a}\mapsto \overline{a}\bmod \langle \overline{-1+2^{e-1}}\rangle$
  is injective. As both groups have the same cardinality $2^{e-2}$, this 
  shows (5).\\
  Assertion (6) is evident for $e\leq 2$. In the remaining situations we know
  that $\textrm{Gal}(K_{2^e}/\QQ)$ is a cyclic group of order $2^{e-2}$.
  This group contains a unique element of order $2$, and evidently this
  element is `complex conjugation'. Moreover, every nontrivial
  subgroup of the Galois group will contain this element. As a consequence,
  the subfield of $K_{2^e}$ corresponding to such a subgroup, is totally real.
  This completes the proof.
\end{proof}

\end{document}
